% =====================================================================
% Figure for lit-review §2.1 — Attack Flow excerpt (Tesla Kubernetes
% cryptojacking, 2018).
%
% Adapted from MITRE/CTID Attack Flow example corpus
% (`Tesla_Kubernetes_Breach.afb`). The figure carries §2.1 P3's claim
% that Attack Flow encodes precondition logic, sequencing, and
% dependency structure between MITRE ATT&CK techniques — i.e. it does
% what intrusion-set membership cannot.
%
% Pedagogical features visible in this excerpt:
%   (1) Condition object (green) — an environmental precondition (the
%       publicly-exposed Kubernetes dashboard) that gates the
%       initial-access action; only the `true' branch is realised.
%   (2) Fan-out from T1133 — two parallel attack paths (credential
%       exfiltration; cryptomining), demonstrating Attack Flow's
%       multi-path representation.
%   (3) AND operator (red) — three-input join in the cryptomining
%       branch; T1496 Resource Hijacking required deploy-container,
%       proxy, and non-standard-port all to succeed. The structural
%       feature a linear kill-chain or intrusion-set cannot express.
%   (4) Infrastructure object (grey cylinder) — a STIX-typed entity
%       (Unlisted mining pool) related to both T1583.004 and T1496
%       via `related-to' edges, illustrating that Attack Flow nodes
%       carry STIX objects beyond techniques themselves. Cylinder
%       shape signals that this is a different kind of node from
%       the action rectangles.
%   (5) Dagger marker (T1552) — the credential-branch path is analyst
%       speculation in the source, not observed activity; the figure
%       preserves provenance fidelity rather than presenting it as
%       fact. Explained in the in-figure legend.
%
% Edge conventions:
%   - Solid arrow (->): `effect' edge — downstream technique cannot
%     execute until upstream completes (Attack Flow's precondition
%     semantics).
%   - Dashed line (no head): `related-to' STIX relationship — no
%     causal ordering implied. Routed along outer columns (x=3.4 and
%     x=13.6), descending below the action grid and turning into
%     Infrastructure's west/east anchors, so neither dashed line
%     crosses any node.
%
% Required packages (loaded in main.tex preamble):
%   \usepackage{tikz}
%   \usetikzlibrary{positioning, arrows.meta, calc, shapes.geometric}
%   \usepackage{xcolor}
% Note: shapes.geometric is required for the cylinder infra style.
%
% Usage:
%   % =====================================================================
% Figure for lit-review §2.1 — Attack Flow excerpt (Tesla Kubernetes
% cryptojacking, 2018).
%
% Adapted from MITRE/CTID Attack Flow example corpus
% (`Tesla_Kubernetes_Breach.afb`). The figure carries §2.1 P3's claim
% that Attack Flow encodes precondition logic, sequencing, and
% dependency structure between MITRE ATT&CK techniques — i.e. it does
% what intrusion-set membership cannot.
%
% Pedagogical features visible in this excerpt:
%   (1) Condition object (green) — an environmental precondition (the
%       publicly-exposed Kubernetes dashboard) that gates the
%       initial-access action; only the `true' branch is realised.
%   (2) Fan-out from T1133 — two parallel attack paths (credential
%       exfiltration; cryptomining), demonstrating Attack Flow's
%       multi-path representation.
%   (3) AND operator (red) — three-input join in the cryptomining
%       branch; T1496 Resource Hijacking required deploy-container,
%       proxy, and non-standard-port all to succeed. The structural
%       feature a linear kill-chain or intrusion-set cannot express.
%   (4) Infrastructure object (grey cylinder) — a STIX-typed entity
%       (Unlisted mining pool) related to both T1583.004 and T1496
%       via `related-to' edges, illustrating that Attack Flow nodes
%       carry STIX objects beyond techniques themselves. Cylinder
%       shape signals that this is a different kind of node from
%       the action rectangles.
%   (5) Dagger marker (T1552) — the credential-branch path is analyst
%       speculation in the source, not observed activity; the figure
%       preserves provenance fidelity rather than presenting it as
%       fact. Explained in the in-figure legend.
%
% Edge conventions:
%   - Solid arrow (->): `effect' edge — downstream technique cannot
%     execute until upstream completes (Attack Flow's precondition
%     semantics).
%   - Dashed line (no head): `related-to' STIX relationship — no
%     causal ordering implied. Routed along outer columns (x=3.4 and
%     x=13.6), descending below the action grid and turning into
%     Infrastructure's west/east anchors, so neither dashed line
%     crosses any node.
%
% Required packages (loaded in main.tex preamble):
%   \usepackage{tikz}
%   \usetikzlibrary{positioning, arrows.meta, calc, shapes.geometric}
%   \usepackage{xcolor}
% Note: shapes.geometric is required for the cylinder infra style.
%
% Usage:
%   % =====================================================================
% Figure for lit-review §2.1 — Attack Flow excerpt (Tesla Kubernetes
% cryptojacking, 2018).
%
% Adapted from MITRE/CTID Attack Flow example corpus
% (`Tesla_Kubernetes_Breach.afb`). The figure carries §2.1 P3's claim
% that Attack Flow encodes precondition logic, sequencing, and
% dependency structure between MITRE ATT&CK techniques — i.e. it does
% what intrusion-set membership cannot.
%
% Pedagogical features visible in this excerpt:
%   (1) Condition object (green) — an environmental precondition (the
%       publicly-exposed Kubernetes dashboard) that gates the
%       initial-access action; only the `true' branch is realised.
%   (2) Fan-out from T1133 — two parallel attack paths (credential
%       exfiltration; cryptomining), demonstrating Attack Flow's
%       multi-path representation.
%   (3) AND operator (red) — three-input join in the cryptomining
%       branch; T1496 Resource Hijacking required deploy-container,
%       proxy, and non-standard-port all to succeed. The structural
%       feature a linear kill-chain or intrusion-set cannot express.
%   (4) Infrastructure object (grey cylinder) — a STIX-typed entity
%       (Unlisted mining pool) related to both T1583.004 and T1496
%       via `related-to' edges, illustrating that Attack Flow nodes
%       carry STIX objects beyond techniques themselves. Cylinder
%       shape signals that this is a different kind of node from
%       the action rectangles.
%   (5) Dagger marker (T1552) — the credential-branch path is analyst
%       speculation in the source, not observed activity; the figure
%       preserves provenance fidelity rather than presenting it as
%       fact. Explained in the in-figure legend.
%
% Edge conventions:
%   - Solid arrow (->): `effect' edge — downstream technique cannot
%     execute until upstream completes (Attack Flow's precondition
%     semantics).
%   - Dashed line (no head): `related-to' STIX relationship — no
%     causal ordering implied. Routed along outer columns (x=3.4 and
%     x=13.6), descending below the action grid and turning into
%     Infrastructure's west/east anchors, so neither dashed line
%     crosses any node.
%
% Required packages (loaded in main.tex preamble):
%   \usepackage{tikz}
%   \usetikzlibrary{positioning, arrows.meta, calc, shapes.geometric}
%   \usepackage{xcolor}
% Note: shapes.geometric is required for the cylinder infra style.
%
% Usage:
%   % =====================================================================
% Figure for lit-review §2.1 — Attack Flow excerpt (Tesla Kubernetes
% cryptojacking, 2018).
%
% Adapted from MITRE/CTID Attack Flow example corpus
% (`Tesla_Kubernetes_Breach.afb`). The figure carries §2.1 P3's claim
% that Attack Flow encodes precondition logic, sequencing, and
% dependency structure between MITRE ATT&CK techniques — i.e. it does
% what intrusion-set membership cannot.
%
% Pedagogical features visible in this excerpt:
%   (1) Condition object (green) — an environmental precondition (the
%       publicly-exposed Kubernetes dashboard) that gates the
%       initial-access action; only the `true' branch is realised.
%   (2) Fan-out from T1133 — two parallel attack paths (credential
%       exfiltration; cryptomining), demonstrating Attack Flow's
%       multi-path representation.
%   (3) AND operator (red) — three-input join in the cryptomining
%       branch; T1496 Resource Hijacking required deploy-container,
%       proxy, and non-standard-port all to succeed. The structural
%       feature a linear kill-chain or intrusion-set cannot express.
%   (4) Infrastructure object (grey cylinder) — a STIX-typed entity
%       (Unlisted mining pool) related to both T1583.004 and T1496
%       via `related-to' edges, illustrating that Attack Flow nodes
%       carry STIX objects beyond techniques themselves. Cylinder
%       shape signals that this is a different kind of node from
%       the action rectangles.
%   (5) Dagger marker (T1552) — the credential-branch path is analyst
%       speculation in the source, not observed activity; the figure
%       preserves provenance fidelity rather than presenting it as
%       fact. Explained in the in-figure legend.
%
% Edge conventions:
%   - Solid arrow (->): `effect' edge — downstream technique cannot
%     execute until upstream completes (Attack Flow's precondition
%     semantics).
%   - Dashed line (no head): `related-to' STIX relationship — no
%     causal ordering implied. Routed along outer columns (x=3.4 and
%     x=13.6), descending below the action grid and turning into
%     Infrastructure's west/east anchors, so neither dashed line
%     crosses any node.
%
% Required packages (loaded in main.tex preamble):
%   \usepackage{tikz}
%   \usetikzlibrary{positioning, arrows.meta, calc, shapes.geometric}
%   \usepackage{xcolor}
% Note: shapes.geometric is required for the cylinder infra style.
%
% Usage:
%   \input{fig_attck_attackflow}
% wrapped in a figure* environment in sections/02_threat_landscape.tex.
%
% Sizing (21 May 2026): natural size approx 16.1 x 7.3 cm. Under the
% current IEEEtran a4paper class, \textwidth approx 18.14 cm, so the
% figure is \input at natural width (NO \resizebox) — it sits inside
% \textwidth with a small margin and is not upscaled. The earlier
% article-class \resizebox{\textwidth}{!}{...} wrapper upscaled the
% artwork approx 13% (inflating its height) and has been removed.
% Vertical spacing was tightened (this revision) to reduce footprint;
% horizontal layout is unchanged. If a future class makes it overflow
% \textwidth, re-wrap the \input in \resizebox{\textwidth}{!}{...}.
% =====================================================================

\begin{tikzpicture}[
  >={Latex[length=2mm,width=1.5mm]},
  every node/.style={font=\scriptsize, inner sep=0pt},
  /tikz/every node/.append style={
    execute at begin node={\hyphenpenalty=10000\exhyphenpenalty=10000}
  },
  action/.style={
    rectangle, rounded corners=2pt,
    draw=blue!55!black, line width=0.45pt,
    fill=blue!5,
    text width=22mm, align=center,
    minimum height=10mm, inner sep=1mm
  },
  condition/.style={
    rectangle, rounded corners=3pt,
    draw=green!50!black, line width=0.55pt,
    fill=green!10,
    text width=24mm, align=center,
    minimum height=10mm, inner sep=1mm
  },
  op/.style={
    circle, draw=red!65!black, line width=0.55pt,
    fill=red!15,
    minimum size=10mm, inner sep=0pt,
    font=\footnotesize\bfseries
  },
  infra/.style={
    cylinder, shape border rotate=90, aspect=0.18,
    draw=black!55, line width=0.5pt,
    fill=black!7,
    text width=22mm, align=center,
    minimum height=11mm, inner sep=1mm,
    font=\scriptsize
  },
  effect/.style={->, semithick, black!75},
  trunk/.style={semithick, black!75},
  related/.style={dashed, semithick, black!55}
]

% ---- Nodes ----
% Top credential row (y = 1.7)
\node[condition] (cond)  at (0,    1.7)  {%
  \textbf{Precondition}\\
  Kubernetes dashboard\\
  exposed; no password};

\node[action]    (t1133) at (3.4,  1.7)  {%
  \textbf{T1133}\\
  External Remote Services\\
  \emph{console login}};

\node[action]    (t1552) at (6.8,  1.7)  {%
  \textbf{T1552}$^{\dagger}$\\
  Unsecured Credentials\\
  \emph{view AWS keys}};

\node[action]    (t1078) at (10.2, 1.7)  {%
  \textbf{T1078.004}\\
  Cloud Accounts\\
  \emph{auth.\ to AWS S3}};

\node[action]    (t1530) at (13.6, 1.7)  {%
  \textbf{T1530}\\
  Data from Cloud Storage\\
  \emph{access S3 buckets}};

% Cryptomining fan-in column (tightened: ~1.35--1.45 cm centre-to-centre)
\node[action]    (t1610) at (6.8,  0.40) {%
  \textbf{T1610}\\
  Deploy Container\\
  \emph{new container}};

\node[action]    (t1090) at (6.8, -0.90) {%
  \textbf{T1090}\\
  Proxy\\
  \emph{Cloudflare CDN}};

\node[action]    (t1571) at (6.8, -2.30) {%
  \textbf{T1571}\\
  Non-Standard Port\\
  \emph{non-default port}};

\node[action]    (t1583) at (3.4, -1.60) {%
  \textbf{T1583.004}\\
  Server\\
  \emph{host mining pool}};

\node[op]        (and)   at (10.2,-0.90) {AND};

\node[action]    (t1496) at (13.6,-0.90) {%
  \textbf{T1496}\\
  Resource Hijacking\\
  \emph{run cryptominer}};

\node[infra]     (infra) at (10.2, -3.45) {%
  \textbf{Infrastructure}\\
  Unlisted mining pool};

% ---- Effect edges (solid; Attack Flow precondition semantics) ----
\draw[effect] (cond)  -- node[above, font=\scriptsize, midway] {true} (t1133);
\draw[effect] (t1133) -- (t1552);
\draw[effect] (t1552) -- (t1078);
\draw[effect] (t1078) -- (t1530);

% T1133 fans down to T1610 (initial access branches into cryptomining)
\draw[effect] (t1133.south) |- (t1610.west);

% T1583 trunks east then splits up to T1090 and down to T1571
\coordinate (tsplit) at (5.5, -1.60);
\draw[trunk]  (t1583.east) -- (tsplit);
\draw[effect] (tsplit) |- (t1090.west);
\draw[effect] (tsplit) |- (t1571.west);

% Three-input fan-in to AND (top, left, bottom)
\draw[effect] (t1610.east) -| (and.north);
\draw[effect] (t1090.east) -- (and.west);
\draw[effect] (t1571.east) -| (and.south);

\draw[effect] (and.east)   -- (t1496.west);

% ---- Related-to edges (dashed; STIX relationship, no causal order) ----
% Routed along the outer columns (x=3.4 from T1583, x=13.6 from T1496),
% descending below the action grid before turning in to Infrastructure's
% west/east anchors. Neither dashed line crosses a node.
\draw[related] (t1583.south) |- (infra.west);
\draw[related] (t1496.south) |- (infra.east);

% ---- In-figure legend strip ----
\begin{scope}[shift={(0, -4.5)}, font=\scriptsize]
  \draw[effect] (0.0,0) -- (0.7,0);
  \node[anchor=west] at (0.75,0) {\emph{effect} (precondition)};
  \draw[related] (4.5,0) -- (5.2,0);
  \node[anchor=west] at (5.25,0) {\emph{related-to} (no causal order)};
  \node[anchor=west] at (9.6,0) {$\dagger$ analyst speculation, not observed};
\end{scope}

\end{tikzpicture}

% wrapped in a figure* environment in sections/02_threat_landscape.tex.
%
% Sizing (21 May 2026): natural size approx 16.1 x 7.3 cm. Under the
% current IEEEtran a4paper class, \textwidth approx 18.14 cm, so the
% figure is \input at natural width (NO \resizebox) — it sits inside
% \textwidth with a small margin and is not upscaled. The earlier
% article-class \resizebox{\textwidth}{!}{...} wrapper upscaled the
% artwork approx 13% (inflating its height) and has been removed.
% Vertical spacing was tightened (this revision) to reduce footprint;
% horizontal layout is unchanged. If a future class makes it overflow
% \textwidth, re-wrap the \input in \resizebox{\textwidth}{!}{...}.
% =====================================================================

\begin{tikzpicture}[
  >={Latex[length=2mm,width=1.5mm]},
  every node/.style={font=\scriptsize, inner sep=0pt},
  /tikz/every node/.append style={
    execute at begin node={\hyphenpenalty=10000\exhyphenpenalty=10000}
  },
  action/.style={
    rectangle, rounded corners=2pt,
    draw=blue!55!black, line width=0.45pt,
    fill=blue!5,
    text width=22mm, align=center,
    minimum height=10mm, inner sep=1mm
  },
  condition/.style={
    rectangle, rounded corners=3pt,
    draw=green!50!black, line width=0.55pt,
    fill=green!10,
    text width=24mm, align=center,
    minimum height=10mm, inner sep=1mm
  },
  op/.style={
    circle, draw=red!65!black, line width=0.55pt,
    fill=red!15,
    minimum size=10mm, inner sep=0pt,
    font=\footnotesize\bfseries
  },
  infra/.style={
    cylinder, shape border rotate=90, aspect=0.18,
    draw=black!55, line width=0.5pt,
    fill=black!7,
    text width=22mm, align=center,
    minimum height=11mm, inner sep=1mm,
    font=\scriptsize
  },
  effect/.style={->, semithick, black!75},
  trunk/.style={semithick, black!75},
  related/.style={dashed, semithick, black!55}
]

% ---- Nodes ----
% Top credential row (y = 1.7)
\node[condition] (cond)  at (0,    1.7)  {%
  \textbf{Precondition}\\
  Kubernetes dashboard\\
  exposed; no password};

\node[action]    (t1133) at (3.4,  1.7)  {%
  \textbf{T1133}\\
  External Remote Services\\
  \emph{console login}};

\node[action]    (t1552) at (6.8,  1.7)  {%
  \textbf{T1552}$^{\dagger}$\\
  Unsecured Credentials\\
  \emph{view AWS keys}};

\node[action]    (t1078) at (10.2, 1.7)  {%
  \textbf{T1078.004}\\
  Cloud Accounts\\
  \emph{auth.\ to AWS S3}};

\node[action]    (t1530) at (13.6, 1.7)  {%
  \textbf{T1530}\\
  Data from Cloud Storage\\
  \emph{access S3 buckets}};

% Cryptomining fan-in column (tightened: ~1.35--1.45 cm centre-to-centre)
\node[action]    (t1610) at (6.8,  0.40) {%
  \textbf{T1610}\\
  Deploy Container\\
  \emph{new container}};

\node[action]    (t1090) at (6.8, -0.90) {%
  \textbf{T1090}\\
  Proxy\\
  \emph{Cloudflare CDN}};

\node[action]    (t1571) at (6.8, -2.30) {%
  \textbf{T1571}\\
  Non-Standard Port\\
  \emph{non-default port}};

\node[action]    (t1583) at (3.4, -1.60) {%
  \textbf{T1583.004}\\
  Server\\
  \emph{host mining pool}};

\node[op]        (and)   at (10.2,-0.90) {AND};

\node[action]    (t1496) at (13.6,-0.90) {%
  \textbf{T1496}\\
  Resource Hijacking\\
  \emph{run cryptominer}};

\node[infra]     (infra) at (10.2, -3.45) {%
  \textbf{Infrastructure}\\
  Unlisted mining pool};

% ---- Effect edges (solid; Attack Flow precondition semantics) ----
\draw[effect] (cond)  -- node[above, font=\scriptsize, midway] {true} (t1133);
\draw[effect] (t1133) -- (t1552);
\draw[effect] (t1552) -- (t1078);
\draw[effect] (t1078) -- (t1530);

% T1133 fans down to T1610 (initial access branches into cryptomining)
\draw[effect] (t1133.south) |- (t1610.west);

% T1583 trunks east then splits up to T1090 and down to T1571
\coordinate (tsplit) at (5.5, -1.60);
\draw[trunk]  (t1583.east) -- (tsplit);
\draw[effect] (tsplit) |- (t1090.west);
\draw[effect] (tsplit) |- (t1571.west);

% Three-input fan-in to AND (top, left, bottom)
\draw[effect] (t1610.east) -| (and.north);
\draw[effect] (t1090.east) -- (and.west);
\draw[effect] (t1571.east) -| (and.south);

\draw[effect] (and.east)   -- (t1496.west);

% ---- Related-to edges (dashed; STIX relationship, no causal order) ----
% Routed along the outer columns (x=3.4 from T1583, x=13.6 from T1496),
% descending below the action grid before turning in to Infrastructure's
% west/east anchors. Neither dashed line crosses a node.
\draw[related] (t1583.south) |- (infra.west);
\draw[related] (t1496.south) |- (infra.east);

% ---- In-figure legend strip ----
\begin{scope}[shift={(0, -4.5)}, font=\scriptsize]
  \draw[effect] (0.0,0) -- (0.7,0);
  \node[anchor=west] at (0.75,0) {\emph{effect} (precondition)};
  \draw[related] (4.5,0) -- (5.2,0);
  \node[anchor=west] at (5.25,0) {\emph{related-to} (no causal order)};
  \node[anchor=west] at (9.6,0) {$\dagger$ analyst speculation, not observed};
\end{scope}

\end{tikzpicture}

% wrapped in a figure* environment in sections/02_threat_landscape.tex.
%
% Sizing (21 May 2026): natural size approx 16.1 x 7.3 cm. Under the
% current IEEEtran a4paper class, \textwidth approx 18.14 cm, so the
% figure is \input at natural width (NO \resizebox) — it sits inside
% \textwidth with a small margin and is not upscaled. The earlier
% article-class \resizebox{\textwidth}{!}{...} wrapper upscaled the
% artwork approx 13% (inflating its height) and has been removed.
% Vertical spacing was tightened (this revision) to reduce footprint;
% horizontal layout is unchanged. If a future class makes it overflow
% \textwidth, re-wrap the \input in \resizebox{\textwidth}{!}{...}.
% =====================================================================

\begin{tikzpicture}[
  >={Latex[length=2mm,width=1.5mm]},
  every node/.style={font=\scriptsize, inner sep=0pt},
  /tikz/every node/.append style={
    execute at begin node={\hyphenpenalty=10000\exhyphenpenalty=10000}
  },
  action/.style={
    rectangle, rounded corners=2pt,
    draw=blue!55!black, line width=0.45pt,
    fill=blue!5,
    text width=22mm, align=center,
    minimum height=10mm, inner sep=1mm
  },
  condition/.style={
    rectangle, rounded corners=3pt,
    draw=green!50!black, line width=0.55pt,
    fill=green!10,
    text width=24mm, align=center,
    minimum height=10mm, inner sep=1mm
  },
  op/.style={
    circle, draw=red!65!black, line width=0.55pt,
    fill=red!15,
    minimum size=10mm, inner sep=0pt,
    font=\footnotesize\bfseries
  },
  infra/.style={
    cylinder, shape border rotate=90, aspect=0.18,
    draw=black!55, line width=0.5pt,
    fill=black!7,
    text width=22mm, align=center,
    minimum height=11mm, inner sep=1mm,
    font=\scriptsize
  },
  effect/.style={->, semithick, black!75},
  trunk/.style={semithick, black!75},
  related/.style={dashed, semithick, black!55}
]

% ---- Nodes ----
% Top credential row (y = 1.7)
\node[condition] (cond)  at (0,    1.7)  {%
  \textbf{Precondition}\\
  Kubernetes dashboard\\
  exposed; no password};

\node[action]    (t1133) at (3.4,  1.7)  {%
  \textbf{T1133}\\
  External Remote Services\\
  \emph{console login}};

\node[action]    (t1552) at (6.8,  1.7)  {%
  \textbf{T1552}$^{\dagger}$\\
  Unsecured Credentials\\
  \emph{view AWS keys}};

\node[action]    (t1078) at (10.2, 1.7)  {%
  \textbf{T1078.004}\\
  Cloud Accounts\\
  \emph{auth.\ to AWS S3}};

\node[action]    (t1530) at (13.6, 1.7)  {%
  \textbf{T1530}\\
  Data from Cloud Storage\\
  \emph{access S3 buckets}};

% Cryptomining fan-in column (tightened: ~1.35--1.45 cm centre-to-centre)
\node[action]    (t1610) at (6.8,  0.40) {%
  \textbf{T1610}\\
  Deploy Container\\
  \emph{new container}};

\node[action]    (t1090) at (6.8, -0.90) {%
  \textbf{T1090}\\
  Proxy\\
  \emph{Cloudflare CDN}};

\node[action]    (t1571) at (6.8, -2.30) {%
  \textbf{T1571}\\
  Non-Standard Port\\
  \emph{non-default port}};

\node[action]    (t1583) at (3.4, -1.60) {%
  \textbf{T1583.004}\\
  Server\\
  \emph{host mining pool}};

\node[op]        (and)   at (10.2,-0.90) {AND};

\node[action]    (t1496) at (13.6,-0.90) {%
  \textbf{T1496}\\
  Resource Hijacking\\
  \emph{run cryptominer}};

\node[infra]     (infra) at (10.2, -3.45) {%
  \textbf{Infrastructure}\\
  Unlisted mining pool};

% ---- Effect edges (solid; Attack Flow precondition semantics) ----
\draw[effect] (cond)  -- node[above, font=\scriptsize, midway] {true} (t1133);
\draw[effect] (t1133) -- (t1552);
\draw[effect] (t1552) -- (t1078);
\draw[effect] (t1078) -- (t1530);

% T1133 fans down to T1610 (initial access branches into cryptomining)
\draw[effect] (t1133.south) |- (t1610.west);

% T1583 trunks east then splits up to T1090 and down to T1571
\coordinate (tsplit) at (5.5, -1.60);
\draw[trunk]  (t1583.east) -- (tsplit);
\draw[effect] (tsplit) |- (t1090.west);
\draw[effect] (tsplit) |- (t1571.west);

% Three-input fan-in to AND (top, left, bottom)
\draw[effect] (t1610.east) -| (and.north);
\draw[effect] (t1090.east) -- (and.west);
\draw[effect] (t1571.east) -| (and.south);

\draw[effect] (and.east)   -- (t1496.west);

% ---- Related-to edges (dashed; STIX relationship, no causal order) ----
% Routed along the outer columns (x=3.4 from T1583, x=13.6 from T1496),
% descending below the action grid before turning in to Infrastructure's
% west/east anchors. Neither dashed line crosses a node.
\draw[related] (t1583.south) |- (infra.west);
\draw[related] (t1496.south) |- (infra.east);

% ---- In-figure legend strip ----
\begin{scope}[shift={(0, -4.5)}, font=\scriptsize]
  \draw[effect] (0.0,0) -- (0.7,0);
  \node[anchor=west] at (0.75,0) {\emph{effect} (precondition)};
  \draw[related] (4.5,0) -- (5.2,0);
  \node[anchor=west] at (5.25,0) {\emph{related-to} (no causal order)};
  \node[anchor=west] at (9.6,0) {$\dagger$ analyst speculation, not observed};
\end{scope}

\end{tikzpicture}

% wrapped in a figure* environment in sections/02_threat_landscape.tex.
%
% Sizing (21 May 2026): natural size approx 16.1 x 7.3 cm. Under the
% current IEEEtran a4paper class, \textwidth approx 18.14 cm, so the
% figure is \input at natural width (NO \resizebox) — it sits inside
% \textwidth with a small margin and is not upscaled. The earlier
% article-class \resizebox{\textwidth}{!}{...} wrapper upscaled the
% artwork approx 13% (inflating its height) and has been removed.
% Vertical spacing was tightened (this revision) to reduce footprint;
% horizontal layout is unchanged. If a future class makes it overflow
% \textwidth, re-wrap the \input in \resizebox{\textwidth}{!}{...}.
% =====================================================================

\begin{tikzpicture}[
  >={Latex[length=2mm,width=1.5mm]},
  every node/.style={font=\scriptsize, inner sep=0pt},
  /tikz/every node/.append style={
    execute at begin node={\hyphenpenalty=10000\exhyphenpenalty=10000}
  },
  action/.style={
    rectangle, rounded corners=2pt,
    draw=blue!55!black, line width=0.45pt,
    fill=blue!5,
    text width=22mm, align=center,
    minimum height=10mm, inner sep=1mm
  },
  condition/.style={
    rectangle, rounded corners=3pt,
    draw=green!50!black, line width=0.55pt,
    fill=green!10,
    text width=24mm, align=center,
    minimum height=10mm, inner sep=1mm
  },
  op/.style={
    circle, draw=red!65!black, line width=0.55pt,
    fill=red!15,
    minimum size=10mm, inner sep=0pt,
    font=\footnotesize\bfseries
  },
  infra/.style={
    cylinder, shape border rotate=90, aspect=0.18,
    draw=black!55, line width=0.5pt,
    fill=black!7,
    text width=22mm, align=center,
    minimum height=11mm, inner sep=1mm,
    font=\scriptsize
  },
  effect/.style={->, semithick, black!75},
  trunk/.style={semithick, black!75},
  related/.style={dashed, semithick, black!55}
]

% ---- Nodes ----
% Top credential row (y = 1.7)
\node[condition] (cond)  at (0,    1.7)  {%
  \textbf{Precondition}\\
  Kubernetes dashboard\\
  exposed; no password};

\node[action]    (t1133) at (3.4,  1.7)  {%
  \textbf{T1133}\\
  External Remote Services\\
  \emph{console login}};

\node[action]    (t1552) at (6.8,  1.7)  {%
  \textbf{T1552}$^{\dagger}$\\
  Unsecured Credentials\\
  \emph{view AWS keys}};

\node[action]    (t1078) at (10.2, 1.7)  {%
  \textbf{T1078.004}\\
  Cloud Accounts\\
  \emph{auth.\ to AWS S3}};

\node[action]    (t1530) at (13.6, 1.7)  {%
  \textbf{T1530}\\
  Data from Cloud Storage\\
  \emph{access S3 buckets}};

% Cryptomining fan-in column (tightened: ~1.35--1.45 cm centre-to-centre)
\node[action]    (t1610) at (6.8,  0.40) {%
  \textbf{T1610}\\
  Deploy Container\\
  \emph{new container}};

\node[action]    (t1090) at (6.8, -0.90) {%
  \textbf{T1090}\\
  Proxy\\
  \emph{Cloudflare CDN}};

\node[action]    (t1571) at (6.8, -2.30) {%
  \textbf{T1571}\\
  Non-Standard Port\\
  \emph{non-default port}};

\node[action]    (t1583) at (3.4, -1.60) {%
  \textbf{T1583.004}\\
  Server\\
  \emph{host mining pool}};

\node[op]        (and)   at (10.2,-0.90) {AND};

\node[action]    (t1496) at (13.6,-0.90) {%
  \textbf{T1496}\\
  Resource Hijacking\\
  \emph{run cryptominer}};

\node[infra]     (infra) at (10.2, -3.45) {%
  \textbf{Infrastructure}\\
  Unlisted mining pool};

% ---- Effect edges (solid; Attack Flow precondition semantics) ----
\draw[effect] (cond)  -- node[above, font=\scriptsize, midway] {true} (t1133);
\draw[effect] (t1133) -- (t1552);
\draw[effect] (t1552) -- (t1078);
\draw[effect] (t1078) -- (t1530);

% T1133 fans down to T1610 (initial access branches into cryptomining)
\draw[effect] (t1133.south) |- (t1610.west);

% T1583 trunks east then splits up to T1090 and down to T1571
\coordinate (tsplit) at (5.5, -1.60);
\draw[trunk]  (t1583.east) -- (tsplit);
\draw[effect] (tsplit) |- (t1090.west);
\draw[effect] (tsplit) |- (t1571.west);

% Three-input fan-in to AND (top, left, bottom)
\draw[effect] (t1610.east) -| (and.north);
\draw[effect] (t1090.east) -- (and.west);
\draw[effect] (t1571.east) -| (and.south);

\draw[effect] (and.east)   -- (t1496.west);

% ---- Related-to edges (dashed; STIX relationship, no causal order) ----
% Routed along the outer columns (x=3.4 from T1583, x=13.6 from T1496),
% descending below the action grid before turning in to Infrastructure's
% west/east anchors. Neither dashed line crosses a node.
\draw[related] (t1583.south) |- (infra.west);
\draw[related] (t1496.south) |- (infra.east);

% ---- In-figure legend strip ----
\begin{scope}[shift={(0, -4.5)}, font=\scriptsize]
  \draw[effect] (0.0,0) -- (0.7,0);
  \node[anchor=west] at (0.75,0) {\emph{effect} (precondition)};
  \draw[related] (4.5,0) -- (5.2,0);
  \node[anchor=west] at (5.25,0) {\emph{related-to} (no causal order)};
  \node[anchor=west] at (9.6,0) {$\dagger$ analyst speculation, not observed};
\end{scope}

\end{tikzpicture}
